% Part: incompleteness
% Chapter: theories-computability
% Section: computably-axiomatizable

\documentclass[../../../include/open-logic-section]{subfiles}

\begin{document}

\olfileid{inc}{tcp}{cax}

\olsection{\printtoken{S}{axiomatizable} સિદ્ધાંતો}

સિદ્ધાંત~$\Th{T}$ પાસે
સ્વયંસિદ્ધિઓનો સંગણનીય ગણ~$A$
હોય તો તેને
\emph{!!{axiomatizable}} કહીએ.
($A$ એ~$\Th{T}$નો
સ્વયંસિદ્ધિઓનો ગણ છે એમ
કહેવાનો અર્થ
$T = \Setabs{!A}{A \Proves !A}$.)
પ્રાકૃતિક સંખ્યાઓના કોઈપણ
``વાજબી'' સ્વયંસિદ્ધીકરણમાં
આ ગુણધર્મ હશે. ખાસ કરીને
સાન્ત સંખ્યામાં સ્વયંસિદ્ધિઓ
ધરાવતો દરેક સિદ્ધાંત
!!{axiomatizable} છે.

\begin{lem}
માનો કે $\Th{T}$
!!{axiomatizable} છે.
તો $\Th{T}$
!!{computably
enumerable} છે.
\end{lem}

\begin{proof}
માનો કે~$A$ એ~$\Th{T}$
માટે સ્વયંસિદ્ધિઓનો
સંગણનીય ગણ છે.
$!A \in T$ હોય તેની
પુષ્ટિ માટે સ્વયંસિદ્ધિઓમાંથી
$!A$નું !!a{derivation}
શોધતા રહો.
\sourcecorrection{OLINC-040}{સ્થિર અંગ્રેજી
મૂળ અહીં શોધથી $!A\in T$
છે કે નહીં ``નક્કી''
થઈ જાય એવું કહે છે.
આ શોધ સભ્યતા હોય ત્યારે
સાબિતી શોધી આપે છે, પરંતુ
સભ્યતા ન હોય ત્યારે અટકવી
જરૂરી નથી; લેમ્માનો દાવો
પરિગણનીયતાનો છે, નિર્ણેયતાનો નહીં.}

થોડું જુદી રીતે કહીએ તો
$!A$ એ $\Th{T}$માં છે
ત્યારે અને માત્ર ત્યારે જ
$A$માં $!B_1$, \dots,
$!B_k$ સ્વયંસિદ્ધિઓની
સાન્ત યાદી અને પ્રથમ-ક્રમના
તર્કમાં
$(!B_1 \land \dots \land !B_k) \lif !A$
નું !!a{derivation} છે.
પરંતુ આપણે અગાઉથી જાણીએ
છીએ કે ``એવા \dots છે કે
\dots'' આકારની વ્યાખ્યા
ધરાવતો દરેક ગણ !!{c.e.}
છે, જો બીજું ``\dots''
સંગણનીય હોય.
\end{proof}

\end{document}
